\section*{Capítulo \ref{ch:g9:fractions} --- Frações e potências}

\begin{solution}{exo:g9:fractions:1}
$\dfrac13 + \dfrac15 = \dfrac{5}{15} + \dfrac{3}{15} = \dfrac{8}{15}$.

$\dfrac72 - \dfrac53 = \dfrac{21}{6} - \dfrac{10}{6} = \dfrac{11}{6}$.

$\dfrac{5}{12} + \dfrac38 = \dfrac{10}{24} + \dfrac{9}{24} =
\dfrac{19}{24}$.

$3 - \dfrac45 = \dfrac{15}{5} - \dfrac45 = \dfrac{11}{5}$.
\end{solution}

\begin{solution}{exo:g9:fractions:2}
$\dfrac59 \times \dfrac{27}{35} = \dfrac{5 \times 27}{9 \times 35}
= \dfrac{27}{9} \times \dfrac{5}{35} = 3 \times \dfrac17 = \dfrac37$
(simplificando por $9$ e por $5$).

$\dfrac{8}{15} \div \dfrac45 = \dfrac{8}{15} \times \dfrac54
= \dfrac{8 \times 5}{15 \times 4} = \dfrac{40}{60} = \dfrac23$.

$\dfrac23 \times \dfrac34 \times \dfrac45 = \dfrac{2 \times 3 \times 4}
{3 \times 4 \times 5} = \dfrac25$ (os $3$ e os $4$ se cancelam).
\end{solution}

\begin{solution}{exo:g9:fractions:3}
$2^5 = 32$; \quad $(-2)^4 = 16$; \quad $-2^4 = -16$; \quad
$4^{-2} = \frac{1}{16}$; \quad
$\left(\frac23\right)^3 = \frac{8}{27}$; \quad $7^0 = 1$.
\end{solution}

\begin{solution}{exo:g9:fractions:4}
$5^3 \times 5^6 = 5^9$; \quad $\dfrac{7^9}{7^5} = 7^4$; \quad
$\left(2^3\right)^4 = 2^{12}$; \quad
$\dfrac{3^2 \times 3^7}{3^5} = 3^{9-5} = 3^4$; \quad
$2^5 \times 4 = 2^5 \times 2^2 = 2^7$.
\end{solution}

\begin{solution}{exo:g9:fractions:5}
$45\,000 = 4.5 \times 10^4$; \quad
$0.0072 = 7.2 \times 10^{-3}$; \quad
$302.5 = 3.025 \times 10^2$; \quad
$0.000\,000\,9 = 9 \times 10^{-7}$; \quad
doze milhões $= 1.2 \times 10^7$.
\end{solution}

\begin{solution}{exo:g9:fractions:6}
$(2 \times 10^7) \times (4.5 \times 10^{-3}) = 9 \times 10^{4}$.

$\dfrac{9 \times 10^5}{3 \times 10^{-2}} = 3 \times 10^{5 - (-2)}
= 3 \times 10^7$.

$(5 \times 10^3)^2 = 25 \times 10^6 = 2.5 \times 10^7$.
\end{solution}

\begin{solution}{exo:g9:fractions:7}
A água acrescentada representa
$\dfrac34 - \dfrac35 = \dfrac{15}{20} - \dfrac{12}{20} =
\dfrac{3}{20}$ do reservatório. Então $\frac{3}{20}$ da \omterm{def:g2:measure:units}{capacidade} são
$30$ litros: $\frac{1}{20}$ são $10$ litros, e a \omterm{def:g2:measure:units}{capacidade} é de $200$
litros.
\end{solution}

\begin{solution}{exo:g9:fractions:8}
Depois de Alice, resta $1 - \frac14 = \frac34$ do bolo. Bento come
$\frac13 \times \frac34 = \frac14$, deixando
$\frac34 - \frac14 = \frac12$. Carla come \omterm{def:g3:division:half}{metade} disso, $\frac14$,
deixando $\frac14$ do bolo.
\end{solution}

\begin{solution}{exo:g9:fractions:9}
Tempo $= \dfrac{\text{distância}}{\text{velocidade}} =
\dfrac{1.5 \times 10^{11}}{3 \times 10^8} = 0.5 \times 10^3 = 500$
segundos, ou seja, $8$ minutos e $20$ segundos.
\end{solution}

\begin{solution}{exo:g9:fractions:10}
Escreva os dois números como \omterm{def:g9:fractions:power}{potências} de \omterm{def:g8:powers:def}{expoente} $10$:
\[
2^{30} = \left(2^3\right)^{10} = 8^{10},
\qquad
3^{20} = \left(3^2\right)^{10} = 9^{10}.
\]
Como $8 < 9$, multiplicando dez cópias de cada um: $8^{10} < 9^{10}$,
então $2^{30} < 3^{20}$.
\end{solution}

\begin{solution}{pb:g9:fractions:1}
\textbf{1.} $\frac14 = 0.25$ (para); $\frac13 = 0.333\ldots$ (bloco
$3$); $\frac56 = 0.8333\ldots$ (bloco $3$ depois do $8$);
$\frac{3}{11} = 0.272727\ldots$ (bloco $27$).

\textbf{2.} A cada passo da divisão por $b$, o \omterm{def:g3:division:remainder}{resto} é um dos $b$
valores $0, 1, \dots, b - 1$ (\cref{thm:g6:wholes:division}). Se
aparecer um \omterm{def:g3:division:remainder}{resto} $0$, a divisão para. Caso contrário, depois de no
máximo $b$ passos algum \omterm{def:g3:division:remainder}{resto} tem de aparecer uma segunda vez --- há
mais passos do que \omterm{def:g3:division:remainder}{restos} possíveis --- e, a partir desse instante, a
divisão repete exatamente a sequência de algarismos que produziu depois
da primeira aparição: a escrita decimal entra em ciclo para sempre.
Parar ou repetir: não existe um terceiro comportamento.

\textbf{3.} $\frac17 = 0.142857\,142857\ldots$: o bloco $142857$ tem
comprimento $6$. A questão 2 prometia uma repetição em no máximo $7$
passos --- aqui, os \omterm{def:g3:division:remainder}{restos} $1, 3, 2, 6, 4, 5$ aparecem todos antes de o
$1$ voltar: a cantoria mais longa possível para uma divisão por $7$.

\textbf{4.} $\frac ab$ para exatamente quando alguma
$\frac{a \times k}{b \times k}$ tem \omterm{def:g4:fractions:def}{denominador} $10$, $100$,
$1000$\,\dots\ --- isto é, quando o \omterm{def:g4:fractions:def}{denominador} pode ser multiplicado
até virar uma \omterm{def:g9:fractions:power}{potência} de dez, como $4 \times 25 = 100$ ou
$8 \times 125 = 1\,000$. Para $3$ e $6$ isso é impossível: uma \omterm{def:g9:fractions:power}{potência}
de dez tem \omterm{def:g1:addition:def}{soma} dos algarismos igual a $1$, e todo \omterm{def:g5:division:multiple}{múltiplo} de $3$ tem
\omterm{def:g1:addition:def}{soma} dos algarismos múltipla de $3$ --- então nenhum \omterm{def:g5:division:multiple}{múltiplo} de $3$
(logo, nenhum de $6$) é jamais $10$, $100$, $1\,000$\,\dots\ Para
$11$: dividir $10$, $100$, $1\,000$, $10\,000$ por $11$ deixa sempre um
\omterm{def:g3:division:remainder}{resto} ($10$, $1$, $10$, $1$\,\dots), nunca $0$. Os \omterm{def:g4:fractions:def}{denominadores} que
alcançam uma \omterm{def:g9:fractions:power}{potência} de dez param; todos os outros se repetem.

\textbf{5.} $\frac{9}{40}$: $40 \times 25 = 1\,000$, para:
$\frac{225}{1000} = 0.225$. $\frac{33}{50}$: $50 \times 2 = 100$,
para: $\frac{66}{100} = 0.66$. $\frac{5}{12}$: $12$ é \omterm{def:g5:division:multiple}{múltiplo} de $3$,
repete ($0.41666\ldots$).

\textbf{6.} $10x = 7.777\ldots$, então $10x - x = 7$ (as caudas se
cancelam perfeitamente), $9x = 7$ e $x = \frac79$. Conferindo pela
divisão: $7 \div 9 = 0.777\ldots$

\textbf{7.} O bloco tem dois algarismos, então desloque por $100$:
$100x - x = 72.7272\ldots - 0.7272\ldots = 72$, logo $99x = 72$ e
$x = \frac{72}{99} = \frac{8}{11}$.

\textbf{8.} $1000x = 583.333\ldots$ e $100x = 58.333\ldots$;
subtraindo, $900x = 525$, então $x = \frac{525}{900} = \frac{7}{12}$
(dividindo por $75$). De fato, $\frac{7}{12} = 0.58333\ldots$

\textbf{9.} $10x - x = 9.999\ldots - 0.999\ldots = 9$, então $9x = 9$
e $x = 1$. Não fica ``logo abaixo'' de $1$: $0.999\ldots$ \emph{é} o
número $1$, escrito com uma segunda fantasia. Todo número estritamente
abaixo de $1$ deixa uma folga, e o \cref{pb:g6:decimals:1} mostrou que
uma folga sempre contém outros números --- ao passo que nada cabe entre
$0.999\ldots$ e $1$. Os noves sem fim são o \omterm{def:g3:division:remainder}{resto} de chocolate do
\cref{pb:g6:fractions:1} encolhido abaixo de qualquer quantidade
positiva: não sobra nada.

\textbf{10.} $2.4545\ldots = 2 + \frac{45}{99} = 2 + \frac{5}{11}
= \frac{27}{11}$. E $0.123123\ldots = \frac{123}{999} =
\frac{41}{333}$ (dividindo por $3$).

\textbf{11.} Pelo dicionário: $0.037037\ldots = \frac{37}{999}$. Como
$999 = 27 \times 37$, isso se simplifica em $\frac{1}{27}$: a \omterm{def:g9:fractions:fraction}{fração}
$\frac{1}{27}$ canta ``$037$''. Conferindo pela divisão:
$1 \div 27 = 0.037037\ldots$

\textbf{12.} $0.0272727\ldots = \frac{1}{10} \times 0.272727\ldots
= \frac{1}{10} \times \frac{27}{99} = \frac{27}{990} = \frac{3}{110}$
(dividindo por $9$).

\textbf{13.} Os blocos de \omterm{def:g1:counting:zero}{zeros} crescem sem parar, então nenhum bloco
fixo pode se repetir para sempre: \emph{não} é uma dízima periódica (e
também nunca para). Pela Parte I, toda \omterm{def:g9:fractions:fraction}{fração} para ou se repete ---
então esse número não é \omterm{def:g9:fractions:fraction}{fração} nenhuma: é irracional, fabricado sob
encomenda. A estrela da irracionalidade, $\sqrt2$, é desmascarada no
\cref{ch:g9:sqrt}.

\textbf{14.} $2^{30} = \left(2^{10}\right)^3 \approx (1.024)^3 \times
10^9 \approx 1.07 \times 10^9$: um número pouco acima de $10^9$, logo
com $10$ algarismos. (Exatamente: $2^{30} = 1\,073\,741\,824$.)

\textbf{15.} $0.20262026\ldots = \frac{2026}{9999}$, já irredutível
($9999 = 3 \times 3 \times 11 \times 101$ não compartilha fator nenhum
com $2026 = 2 \times 1013$). O dicionário completo: \emph{as frações
são exatamente os decimais que param ou se repetem} --- os
\omterm{def:g4:fractions:def}{denominadores} amigos das \omterm{def:g9:fractions:power}{potências} de dez param, todos os outros cantam
um bloco; e, reciprocamente, toda cantoria, até o ano da sua formatura,
é uma \omterm{def:g9:fractions:fraction}{fração} disfarçada.
\end{solution}
